% 12-linearity.tex — Chapter 12: Ownership and Linearity
\chapter{Ownership and Linearity}
\label{chap:12-linearity}
\index{ownership}
\index{linearity}
\index{linear types}

\pnum
The linearity checker enforces Lain's ownership model: every value of a linear
type must be consumed exactly once.  It is implemented in
\srcref{src/sema/linearity.h} and runs after both sema phases via
\cname{sema\_check\_function\_linearity(d)}.

% -------------------------------------------------------------------------
\section{Linear type classification}
\label{sec:lin-class}
\index{sema\_type\_is\_linear}
% -------------------------------------------------------------------------

\pnum
A type is linear if:
\begin{enumerate}
\item Its \cname{OwnershipMode} is \cname{MODE\_OWNED} (set by \lain{mov T}), or
\item It is a fixed-length array (\cname{TYPE\_ARRAY}) whose element type is linear, or
\item It is a \cname{TYPE\_SIMPLE} struct or enum that contains at least one
  linear field (computed recursively).
\end{enumerate}

\pnum
Slices, pointers, and dynamic arrays are \emph{not} linear by default: they are
views, not owners.  Only \cname{TYPE\_ARRAY} (which owns its elements) propagates
linearity from element to container.

% -------------------------------------------------------------------------
\section{Linearity table}
\label{sec:lin-table}
\index{LTable}
\index{LEntry}
% -------------------------------------------------------------------------

\pnum
The checker maintains an \cname{LTable} --- a linked list of \cname{LEntry}
nodes, one per linear variable in the current function:

\begin{structdoc}{LEntry}{src/sema/linearity.h:120}
\cname{id}               & \cname{Id *}       & Variable identifier. \\
\cname{state}            & \cname{LState}     & \cname{LSTATE\_UNCONSUMED} or \cname{LSTATE\_CONSUMED}. \\
\cname{must\_consume}    & \cname{bool}       & True if type is strictly linear. \\
\cname{is\_initialized}  & \cname{bool}       & True if definitely initialized. \\
\cname{is\_defer\_consumed}& \cname{bool}     & True if consumed inside a \lain{defer} block. \\
\cname{field\_states}    & \cname{FieldState *} & Per-field tracking for struct types (see below). \\
\cname{defined\_loop\_depth}& \cname{int}     & Loop nesting depth at declaration. \\
\cname{region}           & \cname{Region *}   & Borrow region at declaration. \\
\cname{line}, \cname{col} & \cname{isize}     & Source location for diagnostics. \\
\end{structdoc}

\pnum
The four \cname{LState} values are:

\begin{center}
\begin{tabular}{ll}
\toprule
\textbf{LState} & \textbf{Meaning} \\
\midrule
\cname{LSTATE\_UNCONSUMED}     & Variable is live and available. \\
\cname{LSTATE\_CONSUMED}       & Variable has been consumed (via \lain{mov} or at last use). \\
\cname{LSTATE\_BORROWED\_READ} & Reserved for future read-borrow state. \\
\cname{LSTATE\_BORROWED\_WRITE}& Reserved for future write-borrow state. \\
\bottomrule
\end{tabular}
\end{center}

% -------------------------------------------------------------------------
\section{Field-level tracking}
\label{sec:lin-fields}
\index{field-level consumption}
% -------------------------------------------------------------------------

\pnum
When a struct variable has multiple linear fields, the checker tracks each
field independently via a linked list of \cname{FieldState} nodes.  A field
is consumed when a \lain{mov} expression targets it:
\lain{mov s.field}.

\pnum
When all linear fields of a struct are consumed, the whole-variable entry is
automatically promoted to \cname{LSTATE\_CONSUMED}.  This allows partial
consumption patterns where individual fields are moved out of a struct without
requiring a whole-struct move.

% -------------------------------------------------------------------------
\section{Consumption rules}
\label{sec:lin-rules}
\index{consumption rules}
% -------------------------------------------------------------------------

\pnum
The checker enforces the following rules:

\begin{description}[leftmargin=2.5cm, style=nextline]
\item[\diagcode{E001}] \textit{Use after move.}  An identifier expression in a
  non-move context on a variable with \cname{LSTATE\_CONSUMED} is an error.
\item[\diagcode{E002}] \textit{Double move.}  A \lain{mov} expression on a
  variable already in \cname{LSTATE\_CONSUMED} is an error.
\item[\diagcode{E003}] \textit{Unconsumed linear value.}  After a function body
  completes, any \cname{must\_consume} variable still in \cname{LSTATE\_UNCONSUMED}
  is an error.
\item[\diagcode{E006}] \textit{Move inside loop.}  Consuming a linear variable
  that was declared outside the current loop nesting level is an error; the
  variable could be consumed on the first iteration and accessed on the second.
\end{description}

% -------------------------------------------------------------------------
\section{Borrow interaction}
\label{sec:lin-borrows}
% -------------------------------------------------------------------------

\pnum
The \cname{LTable} carries a \cname{BorrowTable *borrows} pointer.  Before a
variable can be moved (\lain{mov x}), the checker verifies that no active borrow
of \texttt{x} exists in the borrow table.  This integrates ownership checking
with the region-based borrow checker described in \cref{chap:13-borrows}.

% -------------------------------------------------------------------------
\section{Non-Lexical Lifetimes (NLL)}
\label{sec:lin-nll}
\index{NLL}
\index{Non-Lexical Lifetimes}
% -------------------------------------------------------------------------

\pnum
\srcref{src/sema/use\_analysis.h} implements a use-analysis pre-pass that
identifies the \emph{last use} of each linear variable.  A ``last use'' of a
linear variable that appears as a function argument may be treated as an
implicit move, eliminating the need for an explicit \lain{mov} at every
call site.  The pre-pass annotates \cname{EXPR\_CALL} argument expressions so
the linearity checker can treat them as consuming moves.
